% =========================================================
% Proof: Ring Closed Under Intersection
% Source: volume-i/algebras-of-sets/notes/set-operations-and-hierarchy/notes-set-operations-and-hierarchy.tex
% Mode: standard
% =========================================================

\subsection*{Ring Closed Under Intersection}
\label{prf:ring-closed-intersection}

\begin{remark*}[Return]
\hyperref[prop:ring-closed-intersection]{$\leftarrow$ Back to Proposition (Ring Closed Under Intersection) in Notes}
\end{remark*}

\begin{proof}
Let $\mathcal{R}$ be a ring of sets on $X$ and let $A, B \in \mathcal{R}$.
The claim is that $A \cap B \in \mathcal{R}$.

% TODO: proof to be filled.
% Strategy: Apply (R3) to get $A\setminus B\in\mathcal{R}$, then apply (R3) again
% to get $A\setminus(A\setminus B)\in\mathcal{R}$, then show $A\setminus(A\setminus B)=A\cap B$.
\end{proof}

\begin{remark*}[Proof shape]
% TODO

\FlashProof{1. Apply (R3): $A\setminus B\in\mathcal{R}$. 2. Apply (R3) again: $A\setminus(A\setminus B)\in\mathcal{R}$. 3. Element chase: $x\in A\setminus(A\setminus B)$ iff $x\in A$ and $x\notin A\setminus B$ iff $x\in A$ and $x\in B$ iff $x\in A\cap B$.}
\FlashImplies{prop:hierarchy-inclusions}
\end{remark*}

\begin{remark*}[Dependencies]
\begin{itemize}
  \item \hyperref[def:ring-of-sets]{Definition of ring of sets}: axioms (R1)--(R3)
  \item \hyperref[def:set-operations]{Definition of set operations}: $\cap$ and $\setminus$
\end{itemize}

\FlashUses{def:ring-of-sets, def:set-operations}
\end{remark*}
